\chapter{Metodologia}
\label{cap:metodologia}

\section{Considerações Iniciais}

Este capítulo descreve a metodologia adotada para comparar as arquiteturas monolítica e de
microsserviços. São apresentados a classificação da pesquisa, as etapas do trabalho, os
requisitos implementados nas duas versões da aplicação, as tecnologias utilizadas, a forma como
foi garantida a equivalência funcional entre as versões, o ambiente de testes, os cenários e
perfis de carga, as métricas associadas a cada critério de comparação, o procedimento de
execução e análise dos dados e, por fim, as ameaças à validade dos resultados.

\section{Classificação da Pesquisa}

Quanto à natureza, esta pesquisa é aplicada, pois busca gerar conhecimento para a solução de um
problema prático: a escolha da arquitetura de software mais adequada para um sistema. Quanto à
abordagem, é quali-quantitativa, já que combina medições numéricas de desempenho, de uso de
recursos e de código-fonte com a análise qualitativa do esforço de desenvolvimento e de
implantação. Quanto aos objetivos, classifica-se como exploratória e descritiva: exploratória
por proporcionar maior familiaridade com o problema por meio do levantamento bibliográfico e do
estado da arte, e descritiva por registrar e comparar as características de cada arquitetura
\cite{gil2017}.

Quanto aos procedimentos técnicos, a pesquisa é bibliográfica, na construção do referencial
teórico e do estado da arte, e experimental, na comparação entre as arquiteturas
\cite{gil2017}. A parte experimental segue os princípios de um experimento controlado em
engenharia de software, conforme \citeonline{wohlin2012}: a arquitetura é o fator estudado, com
dois tratamentos (monolítica e de microsserviços), enquanto as demais variáveis, como
funcionalidades, linguagem, \textit{framework}, banco de dados, máquina, recursos disponíveis e
carga aplicada, são mantidas constantes. As variáveis dependentes são as métricas descritas na
Seção~\ref{sec:metricas}.

\section{Etapas do Trabalho}

O trabalho foi organizado nas seguintes etapas:

\begin{enumerate}
    \item \textbf{Revisão bibliográfica:} estudo dos conceitos de arquitetura de software, das
    arquiteturas monolítica e de microsserviços, das formas de comunicação entre componentes e
    dos atributos de qualidade, apresentados no Capítulo~\ref{cap:referencial};
    \item \textbf{Levantamento do estado da arte:} análise dos trabalhos que compararam as duas
    arquiteturas, apresentada no Capítulo~\ref{cap:estado-arte}, que orientou a escolha dos
    critérios e das métricas;
    \item \textbf{Definição dos requisitos:} especificação das funcionalidades comuns às duas
    versões da aplicação de \textit{e-commerce};
    \item \textbf{Desenvolvimento da versão em microsserviços:} implementação dos serviços, do
    \textit{API Gateway} e da comunicação síncrona e assíncrona entre eles;
    \item \textbf{Desenvolvimento da versão monolítica:} implementação das mesmas regras de
    negócio em uma única aplicação, organizada em módulos;
    \item \textbf{Verificação da equivalência funcional:} execução de testes automatizados que
    confirmam que as duas versões oferecem as mesmas rotas e produzem as mesmas respostas;
    \item \textbf{Preparação do ambiente de testes:} conteinerização das aplicações, definição
    dos limites de recursos, simulação do serviço de pagamento e criação dos cenários de carga;
    \item \textbf{Execução dos experimentos e coleta das métricas:} realização dos testes de
    carga, de estresse e de escalabilidade e coleta das métricas de código e de implantação;
    \item \textbf{Análise e discussão dos resultados:} tratamento estatístico dos dados,
    comparação entre as arquiteturas e confronto com os trabalhos do estado da arte.
\end{enumerate}

\section{Requisitos das Aplicações}

As duas versões implementam o mesmo conjunto de requisitos funcionais, apresentado na
Tabela~\ref{tab:requisitos}. Os requisitos foram definidos de modo a reproduzir as operações
típicas de uma loja virtual e a exercitar tanto a comunicação síncrona quanto a assíncrona,
além da integração com serviços externos de pagamento e de armazenamento de arquivos.

\begin{table}[htbp]
\centering
\caption{Requisitos funcionais das aplicações}
\label{tab:requisitos}
\small
\begin{tabular}{|>{\raggedright\arraybackslash}p{1.5cm}|>{\raggedright\arraybackslash}p{13.1cm}|}
\hline
\textbf{Código} & \textbf{Descrição} \\ \hline
RF01 & Cadastrar usuários e autenticá-los por meio de \textit{token} JWT, com os perfis de cliente e de administrador \\ \hline
RF02 & Criar automaticamente o participante, com os dados pessoais do cliente, a partir do cadastro do usuário \\ \hline
RF03 & Consultar participantes, sendo a listagem completa restrita ao administrador \\ \hline
RF04 & Cadastrar, alterar e excluir produtos, com imagem, variações de tamanho e cor e controle de estoque (administrador) \\ \hline
RF05 & Consultar o catálogo de produtos sem necessidade de autenticação: todos, apenas os com estoque, por identificador e por lista de identificadores \\ \hline
RF06 & Personalizar a loja, com cores e logotipo (administrador), e consultar essa configuração publicamente \\ \hline
RF07 & Configurar o \textit{token} de acesso ao Mercado Pago da loja (administrador) \\ \hline
RF08 & Realizar pedidos, validando o cliente, a existência dos produtos e variações e o estoque disponível \\ \hline
RF09 & Processar o pagamento do pedido por PIX, cartão de crédito ou cartão de débito por meio do Mercado Pago \\ \hline
RF10 & Receber as notificações do Mercado Pago, atualizar o status do pagamento do pedido e baixar o estoque quando o pagamento for aprovado \\ \hline
\end{tabular}
\par\vspace{4pt}{\footnotesize Fonte: Elaborada pelo autor (2026).}
\end{table}

Além dos requisitos funcionais, foram definidos requisitos não funcionais comuns às duas
versões. O principal deles é que ambas exponham exatamente as mesmas rotas, com os mesmos
formatos de requisição e de resposta em JSON e os mesmos códigos de \textit{status} HTTP, de
modo que um mesmo cliente, ou um mesmo \textit{script} de teste, possa ser utilizado com
qualquer uma delas. Também foram definidos a autenticação sem estado por JWT, o armazenamento
das imagens em um serviço de objetos compatível com o protocolo S3 e a execução de todos os
componentes em contêineres Docker. As diferenças entre as versões ficam restritas à
organização interna: na versão de microsserviços, cada serviço possui seu próprio banco de
dados e se comunica com os demais por HTTP ou por mensagens; na versão monolítica, há um único
banco de dados e os módulos se comunicam por chamadas diretas e eventos internos. Essas
diferenças são detalhadas no Capítulo~\ref{cap:desenvolvimento}.

\section{Tecnologias Utilizadas}

As duas versões foram desenvolvidas com a mesma pilha tecnológica, de modo que as diferenças
observadas possam ser atribuídas à arquitetura e não às ferramentas. A
Tabela~\ref{tab:tecnologias} apresenta as tecnologias utilizadas e em qual versão cada uma é
empregada.

\begin{table}[htbp]
\centering
\caption{Tecnologias utilizadas no desenvolvimento e nos testes}
\label{tab:tecnologias}
\small
\begin{tabular}{|>{\raggedright\arraybackslash}p{4.4cm}|>{\raggedright\arraybackslash}p{7.4cm}|>{\raggedright\arraybackslash}p{2.8cm}|}
\hline
\textbf{Tecnologia} & \textbf{Finalidade} & \textbf{Versão em que é usada} \\ \hline
Java 21 & Linguagem de programação & Ambas \\ \hline
Spring Boot 3.5.13 &\textit{Framework} das aplicações (Web, segurança, acesso a dados) & Ambas \\ \hline
Spring Security e JWT & Autenticação e autorização sem estado & Ambas \\ \hline
Spring Data JPA e Hibernate & Persistência de dados & Ambas \\ \hline
PostgreSQL 17.4 & Banco de dados relacional & Ambas \\ \hline
MinIO & Armazenamento das imagens (compatível com S3) & Ambas \\ \hline
SDK Java do Mercado Pago 2.8.0 & Integração com o serviço de pagamentos & Ambas \\ \hline
MapStruct e Lombok & Conversão entre objetos e redução de código repetitivo & Ambas \\ \hline
JUnit e Mockito & Testes automatizados & Ambas \\ \hline
Spring Cloud Gateway & Ponto único de entrada (\textit{API Gateway}) & Microsserviços \\ \hline
Spring Cloud OpenFeign & Comunicação síncrona entre serviços (HTTP) & Microsserviços \\ \hline
Apache Kafka (Confluent 7.7.7) & Comunicação assíncrona entre serviços & Microsserviços \\ \hline
Spring Application Events & Comunicação entre módulos por eventos internos & Monolítica \\ \hline
Docker e Docker Compose & Conteinerização e orquestração do ambiente & Ambas \\ \hline
Grafana k6 0.54 & Geração de carga e medição de desempenho \cite{grafana-k6} & Testes \\ \hline
WireMock 3.9 & Simulação da API do Mercado Pago \cite{wiremock} & Testes \\ \hline
Nginx 1.27 & Balanceador de carga nos testes de escalabilidade horizontal & Testes \\ \hline
\end{tabular}
\par\vspace{4pt}{\footnotesize Fonte: Elaborada pelo autor (2026).}
\end{table}

\section{Verificação da Equivalência Funcional}

Para que a comparação seja válida, as duas versões precisam se comportar da mesma forma para
quem as utiliza. Por isso, além dos testes unitários de cada versão, foram criados
\textit{scripts} de teste funcional que acessam a aplicação somente pela sua interface
externa, por meio de requisições HTTP, e que são executados sem alteração contra as duas
versões. Esses \textit{scripts} verificam cada rota da aplicação, incluindo os casos de erro,
como requisições sem autenticação, dados inválidos, estoque insuficiente e falhas na
comunicação com o Mercado Pago, além de cenários completos de uso, como o cadastro de um
cliente seguido da realização de um pedido. Para cada verificação, são comparados o código de
\textit{status} e o corpo da resposta.

O fluxo assíncrono de pagamento também é verificado: após a notificação de pagamento aprovado,
confere-se se o status do pedido foi atualizado e se o estoque dos produtos foi baixado. Sempre
que uma regra de negócio é alterada, a alteração é aplicada nas duas versões e os testes são
executados novamente. Os resultados dessa verificação são apresentados no
Capítulo~\ref{cap:desenvolvimento}.

\section{Ambiente de Testes}

Os experimentos são executados em um único computador, cuja configuração é apresentada na
Tabela~\ref{tab:maquina}. Todos os componentes, inclusive o gerador de carga, são executados
em contêineres Docker conectados à mesma rede virtual, o que torna o ambiente reproduzível: o
mesmo conjunto de arquivos de configuração e de \textit{scripts} cria o ambiente do zero, com
os bancos de dados vazios, a cada execução.

\begin{table}[htbp]
\centering
\caption{Configuração da máquina utilizada nos experimentos}
\label{tab:maquina}
\small
\begin{tabular}{|>{\raggedright\arraybackslash}p{4.4cm}|>{\raggedright\arraybackslash}p{10.2cm}|}
\hline
\textbf{Componente} & \textbf{Especificação} \\ \hline
Processador & AMD Ryzen 5 4600G, 6 núcleos e 12 \textit{threads} \\ \hline
Memória & 32 GB \\ \hline
Armazenamento & SSD SATA de 960 GB \\ \hline
Sistema operacional & Windows 11 Pro \\ \hline
Contêineres & Docker Desktop 29.9 sobre WSL 2, com 12 CPUs e 15 GB de memória disponíveis \\ \hline
Máquina virtual Java & Eclipse Temurin 21 (JRE), igual em todas as imagens \\ \hline
\end{tabular}
\par\vspace{4pt}{\footnotesize Fonte: Elaborada pelo autor (2026).}
\end{table}

A Figura~\ref{fig:ambiente-testes} ilustra a organização do ambiente. O gerador de carga envia
as requisições diretamente ao ponto de entrada da aplicação sob teste, que é o próprio monólito
ou o \textit{API Gateway} dos microsserviços. A aplicação utiliza o PostgreSQL, o MinIO e,
apenas na versão de microsserviços, o Kafka. Durante a medição, um coletor registra
periodicamente o consumo de CPU e de memória de cada contêiner. As duas arquiteturas nunca são
executadas ao mesmo tempo, e os ambientes de desenvolvimento e demais programas da máquina
permanecem encerrados durante os experimentos.

\begin{figure}[htbp]
\centering
\caption{Organização do ambiente de testes}
\label{fig:ambiente-testes}
\resizebox{\textwidth}{!}{%
\begin{tikzpicture}[
    caixa/.style={draw, rounded corners=2pt, minimum width=3.4cm, minimum height=1.0cm, fill=gray!10, font=\small, align=center},
    bd/.style={draw, cylinder, shape border rotate=90, aspect=0.25, minimum width=2.2cm, minimum height=1.1cm, fill=gray!25, font=\small},
    seta/.style={-{Stealth[length=2.2mm]}, thick}
]
\node[caixa, fill=white] (k6) at (0,0) {Gerador de carga\\(k6)};
\node[caixa, fill=gray!30, minimum width=6.4cm, minimum height=1.8cm] (app) at (6.2,0) {Aplicação sob teste\\monólito \textit{ou}\\\textit{API Gateway} e microsserviços};
\node[caixa] (mp) at (13.2,0) {Simulador do\\Mercado Pago (WireMock)};
\node[bd, minimum width=2.4cm] (pg) at (3.6,-3.0) {PostgreSQL};
\node[caixa, minimum width=2.2cm] (minio) at (6.2,-3.0) {MinIO};
\node[caixa] (kafka) at (9.4,-3.0) {Kafka e Zookeeper\\(só microsserviços)};
\node[caixa, fill=white] (coletor) at (0,-3.0) {Coletor de recursos\\(\textit{docker stats})};
\draw[seta] (k6) -- node[above, font=\footnotesize] {HTTP} (app);
\draw[seta] (app) -- node[above, font=\footnotesize] {HTTPS} (mp);
\draw[seta] (app.south -| pg) -- (pg.north);
\draw[seta] (app.south -| minio) -- (minio.north);
\draw[seta, dashed] (app.south -| {(8.6,0)}) -- ({(8.6,0)} |- kafka.north);
\draw[seta, dotted] (coletor.north) |- ([yshift=0.3cm]app.south west);
\node[draw, thick, dashed, rounded corners=4pt, inner sep=10pt, fit=(k6)(mp)(coletor)(kafka), label={[font=\footnotesize]below:Rede virtual do Docker em uma única máquina}] {};
\end{tikzpicture}%
}
\par\vspace{4pt}{\footnotesize Fonte: Elaborada pelo autor (2026).}
\end{figure}

\subsection{Limites de Recursos}

Para que a comparação seja justa, a camada de aplicação recebe o mesmo orçamento de recursos
nas duas arquiteturas: 4 CPUs e 4 GB de memória no total. No monólito, todo esse orçamento é
atribuído ao único contêiner da aplicação; nos microsserviços, ele é dividido entre os sete
contêineres, proporcionalmente à participação esperada de cada serviço nas operações. Os
componentes de infraestrutura recebem os mesmos limites nas duas arquiteturas, e o Kafka e o
Zookeeper, existentes apenas na versão de microsserviços, recebem limites próprios. A
Tabela~\ref{tab:limites} apresenta os limites utilizados. Todas as aplicações são executadas
com a mesma imagem base e com a mesma configuração da máquina virtual Java, que limita o
tamanho do \textit{heap} a 60\% da memória do contêiner, e o registro dos comandos SQL
executados é desativado nas duas versões.

\begin{table}[htbp]
\centering
\caption{Limites de recursos dos contêineres}
\label{tab:limites}
\small
\begin{tabular}{|>{\raggedright\arraybackslash}p{5.6cm}|>{\centering\arraybackslash}p{1.6cm}|>{\centering\arraybackslash}p{2.8cm}|>{\raggedright\arraybackslash}p{3.4cm}|}
\hline
\textbf{Contêiner} & \textbf{CPUs} & \textbf{Memória} & \textbf{Arquitetura} \\ \hline
Monólito & 4,00 & 4.096 MB & Monolítica \\ \hline
\textit{API Gateway} & 0,75 & 640 MB & Microsserviços \\ \hline
Auth & 0,50 & 512 MB & Microsserviços \\ \hline
Participante & 0,50 & 512 MB & Microsserviços \\ \hline
Produto & 0,75 & 640 MB & Microsserviços \\ \hline
Loja & 0,50 & 512 MB & Microsserviços \\ \hline
Pedido & 0,50 & 640 MB & Microsserviços \\ \hline
Pagamento & 0,50 & 512 MB & Microsserviços \\ \hline
\textbf{Total da aplicação} & \textbf{4,00} & \textbf{$\approx$ 4.000 MB} & \textbf{Ambas} \\ \hline
PostgreSQL & 2,00 & 1.024 MB & Ambas \\ \hline
MinIO & 0,50 & 512 MB & Ambas \\ \hline
Simulador do Mercado Pago & 1,00 & 512 MB & Ambas \\ \hline
Kafka & 1,00 & 1.024 MB & Microsserviços \\ \hline
Zookeeper & 0,50 & 512 MB & Microsserviços \\ \hline
Gerador de carga (k6) & 2,00 & 1.024 MB & Ambas \\ \hline
Balanceador (Nginx)\textsuperscript{*} & 1,00 & 256 MB & Ambas \\ \hline
\end{tabular}
\par\vspace{4pt}{\footnotesize \textsuperscript{*}Utilizado apenas nos testes de escalabilidade horizontal.\\Fonte: Elaborada pelo autor (2026).}
\end{table}

\subsection{Simulação do Serviço de Pagamento}

O processamento de pagamentos depende da API do Mercado Pago, um serviço externo cujo tempo de
resposta varia com fatores alheios ao experimento, como a conexão com a Internet e a carga do
próprio serviço, e cujo uso intensivo durante os testes não seria adequado. Por isso, a API foi
substituída por um simulador construído com o WireMock \cite{wiremock}, que atende pelo mesmo
endereço da API real dentro da rede do Docker, utilizando um certificado digital criado
especificamente para os testes. Dessa forma, o código de pagamento das aplicações permanece
inalterado. O simulador responde a todas as chamadas com um atraso fixo de 100~ms,
representando o tempo de uma chamada a um serviço externo, cria os pagamentos por PIX como
pendentes e os pagamentos por cartão como aprovados e informa, nas consultas, que o pagamento
foi aprovado. Como o atraso é o mesmo nas duas arquiteturas, ele não favorece nenhuma delas.

\subsection{Massa de Dados}

Antes de cada execução, os bancos de dados são recriados e populados, por meio da própria
interface da aplicação, com a mesma massa de dados: um administrador, cem clientes, vinte
produtos com estoque suficiente para não se esgotar durante os testes, sendo quatro deles com
três variações cada, e a configuração do Mercado Pago da loja. Assim, as duas arquiteturas
partem sempre do mesmo estado inicial.

\section{Cenários e Perfis de Carga}

Os testes de desempenho são realizados com o Grafana k6, ferramenta de código aberto para
testes de carga em que os cenários são escritos em JavaScript e que registra o tempo de cada
requisição \cite{grafana-k6}. Foram definidos três cenários, apresentados na
Tabela~\ref{tab:cenarios}. No cenário de compra, cada pedido percorre os módulos de pedido,
produto, participante, pagamento, loja e autenticação, além do simulador do Mercado Pago, e a
notificação de pagamento desencadeia a atualização do pedido e a baixa do estoque, que nos
microsserviços ocorre de forma assíncrona por meio do Kafka.

\begin{table}[htbp]
\centering
\caption{Cenários de teste}
\label{tab:cenarios}
\small
\begin{tabular}{|>{\raggedright\arraybackslash}p{2.4cm}|>{\raggedright\arraybackslash}p{12.2cm}|}
\hline
\textbf{Cenário} & \textbf{Operações realizadas em cada iteração} \\ \hline
Catálogo & Lista os produtos com estoque, consulta um produto escolhido aleatoriamente e consulta a configuração da loja, sem autenticação \\ \hline
Compra & Um cliente autenticado realiza um pedido com dois produtos escolhidos aleatoriamente, pago por PIX, e em seguida é enviada a notificação de pagamento aprovado \\ \hline
Misto & Executa os dois cenários anteriores em paralelo, na proporção de 80\% de iterações de catálogo para 20\% de compra, aproximando o comportamento de uma loja virtual, em que a maior parte dos acessos é de navegação \\ \hline
\end{tabular}
\par\vspace{4pt}{\footnotesize Fonte: Elaborada pelo autor (2026).}
\end{table}

Os cenários são executados segundo três perfis de carga. No teste de fumaça, um único usuário
virtual executa poucas iterações, apenas para confirmar que o ambiente está funcionando. No
teste de carga, o k6 inicia novas iterações a uma taxa constante durante três minutos,
independentemente do tempo de resposta da aplicação; esse modelo, chamado de aberto, evita que
uma aplicação mais lenta receba menos requisições e mascare a sua lentidão. As taxas
utilizadas são as mesmas para as duas arquiteturas e são definidas a partir de um teste
preliminar, de modo a representar cargas baixa, média e alta. No teste de estresse, a taxa é
aumentada em oito degraus iguais até um valor máximo, com o objetivo de encontrar o ponto de
saturação de cada arquitetura. Cada degrau dura um minuto no teste preliminar e trinta segundos
nas execuções repetidas, o que reduz a duração total dos testes sem alterar o número de degraus.

A máquina virtual Java compila e otimiza o código durante a execução, de modo que as primeiras
requisições são mais lentas do que as seguintes. Por isso, antes de cada medição, é executado
um aquecimento de um minuto com metade da taxa do teste, cujos resultados são descartados.

\section{Métricas e Critérios de Comparação}
\label{sec:metricas}

Os cinco atributos de qualidade descritos no Capítulo~\ref{cap:referencial} são utilizados
como critérios de comparação. A Tabela~\ref{tab:metricas} relaciona cada critério às métricas
utilizadas e aos instrumentos de coleta.

\begin{table}[htbp]
\centering
\caption{Critérios de comparação, métricas e instrumentos de coleta}
\label{tab:metricas}
\small
\begin{tabular}{|>{\raggedright\arraybackslash}p{3.1cm}|>{\raggedright\arraybackslash}p{7.5cm}|>{\raggedright\arraybackslash}p{4.0cm}|}
\hline
\textbf{Critério} & \textbf{Métricas} & \textbf{Instrumento} \\ \hline
Desempenho & Tempo de resposta (média, mediana e percentis 90, 95 e 99), vazão e taxa de erros, por cenário e por rota; uso de CPU e de memória sob carga & k6 e \textit{docker stats} \\ \hline
Escalabilidade & Vazão máxima sustentada; variação da vazão e do tempo de resposta com o aumento dos recursos e do número de instâncias; eficiência da escala & k6 e \textit{docker stats} \\ \hline
Manutenibilidade & Linhas de código, número de classes, complexidade ciclomática, acoplamento entre objetos (CBO) e falta de coesão em métodos (LCOM); impacto de alterações & Ferramenta CK \cite{aniche-ck} e análise dos repositórios \\ \hline
Complexidade de desenvolvimento & Número de projetos, de arquivos de configuração e de dependências; quantidade de código dedicada à comunicação entre componentes; tempo de compilação & Análise dos repositórios \\ \hline
Implantação & Número de imagens e contêineres; tamanho das imagens; tempo de construção e de inicialização; consumo de recursos em repouso & Docker e \textit{scripts} do ambiente \\ \hline
\end{tabular}
\par\vspace{4pt}{\footnotesize Fonte: Elaborada pelo autor (2026).}
\end{table}

No critério de desempenho, o tempo de resposta é analisado principalmente por meio dos
percentis, que revelam o comportamento das requisições mais lentas, ocultado pela média
\cite{jain1991}. Os resultados são registrados por rota e por cenário, permitindo identificar
quais operações são mais afetadas pela comunicação entre serviços.

Na escalabilidade, considera-se como vazão máxima sustentada a maior taxa, obtida no teste de
estresse, em que a taxa de erros permanece abaixo de 1\% e o percentil 95 do tempo de resposta
permanece abaixo de um segundo. A escalabilidade vertical é avaliada repetindo o teste de
estresse com orçamentos de 2, 4 e 8 CPUs para a camada de aplicação, mantendo a memória em
4 GB, já que reduzi-la na mesma proporção deixaria cada serviço abaixo da memória mínima
necessária para executar uma aplicação Spring Boot. A escalabilidade horizontal é avaliada
acrescentando instâncias atrás de um balanceador de carga Nginx, utilizado da mesma forma nas
duas arquiteturas: no monólito, a aplicação inteira é replicada, correspondendo ao eixo X do
cubo de escalabilidade; nos microsserviços, são replicados apenas os serviços mais demandados,
correspondendo à combinação dos eixos X e Y \cite{abbott2015}. Como o balanceador acrescenta
um intermediário às requisições, a referência para essa comparação é a execução com uma única
instância de cada componente atrás do mesmo balanceador. A eficiência da escala é calculada
pela razão entre o ganho de vazão e o aumento de recursos.

Na manutenibilidade, as métricas de código são extraídas com a ferramenta CK
\cite{aniche-ck}, que calcula as métricas propostas por \citeonline{chidamber1994} e a
complexidade ciclomática de \citeonline{mccabe1976} para código Java. Como complemento, é
analisado o impacto de duas alterações hipotéticas nas duas versões, a inclusão de um novo
atributo no produto que precise ser exibido no pedido e a inclusão de um novo meio de
pagamento, contando quantos arquivos, módulos ou serviços e implantações seriam afetados.

A complexidade de desenvolvimento e a implantação são avaliadas a partir dos repositórios e do
próprio ambiente de testes, contabilizando os artefatos necessários para construir e executar
cada versão. O tempo de inicialização é medido desde o comando que sobe o ambiente até que
todas as aplicações estejam prontas para receber requisições, e o consumo de recursos em
repouso é medido com as aplicações no ar e sem carga.

\section{Procedimento de Execução e Análise dos Dados}

Cada combinação de arquitetura, cenário e perfil de carga é executada cinco vezes. A ordem das
execuções é alternada entre as arquiteturas, de modo que eventuais variações da máquina ao
longo do tempo afetem as duas igualmente. Cada execução segue os passos abaixo, automatizados
por \textit{scripts}:

\begin{enumerate}
    \item recriação do ambiente, com os bancos de dados vazios, e registro do tempo de
    inicialização;
    \item cadastro da massa de dados;
    \item aquecimento da máquina virtual Java, sem registro dos resultados;
    \item execução do cenário, com coleta simultânea das métricas do k6 e do uso de CPU e de
    memória dos contêineres;
    \item registro dos resultados, dos \textit{logs} das aplicações e das informações do
    ambiente, como as versões do código e os limites de recursos utilizados;
    \item remoção do ambiente.
\end{enumerate}

Para cada métrica, são calculados a média, o desvio padrão e o coeficiente de variação das
cinco repetições, além do intervalo de confiança de 95\%, conforme recomendado por
\citeonline{jain1991} para a comparação de sistemas. As arquiteturas são comparadas pela
diferença percentual entre as médias, e uma diferença é considerada significativa apenas
quando os intervalos de confiança das duas arquiteturas não se sobrepõem. Os resultados são
apresentados em tabelas e gráficos no Capítulo~\ref{cap:resultados}.

\section{Ameaças à Validade}

\citeonline{wohlin2012} recomendam que os experimentos em engenharia de software discutam os
fatores que podem comprometer a validade de seus resultados. Neste trabalho, foram
identificadas as seguintes ameaças:

\begin{itemize}
    \item \textbf{Execução em uma única máquina:} a comunicação entre os contêineres ocorre em
    uma rede virtual local, com latência muito menor que a de uma rede física entre
    servidores. Assim, o custo da comunicação entre os microsserviços tende a ser
    subestimado em relação a um ambiente distribuído real;
    \item \textbf{Virtualização:} o Docker Desktop executa os contêineres em uma máquina
    virtual sobre o Windows, o que acrescenta uma camada que não existiria em um servidor
    Linux. Como essa camada é a mesma para as duas arquiteturas, o efeito sobre a comparação
    tende a ser reduzido;
    \item \textbf{Simulação do pagamento:} o atraso fixo do simulador não reproduz a
    variabilidade do serviço real do Mercado Pago;
    \item \textbf{Configurações padrão:} as aplicações utilizam as configurações padrão do Spring
    Boot, como o \textit{pool} de dez conexões com o banco de dados por aplicação. Como o
    monólito tem um único \textit{pool} e cada microsserviço tem o seu, essas configurações
    podem limitar as arquiteturas de formas diferentes. Elas foram mantidas iguais por
    aplicação para não favorecer nenhuma das versões, e os seus efeitos são discutidos nos
    resultados;
    \item \textbf{Divisão dos recursos:} a forma como o orçamento é dividido entre os
    microsserviços influencia o desempenho dessa arquitetura. Por isso, a divisão é
    registrada junto com os resultados e avaliada nos testes de escalabilidade;
    \item \textbf{Autoria das implementações:} as duas versões foram desenvolvidas pelo mesmo
    autor, o que pode favorecer, ainda que involuntariamente, uma delas. Essa ameaça é
    reduzida pelo compartilhamento das regras de negócio e pela verificação automatizada da
    equivalência funcional;
    \item \textbf{Generalização:} os resultados se referem a uma aplicação de porte pequeno a
    médio, com uma equipe de desenvolvimento de uma pessoa, e não podem ser generalizados
    diretamente para sistemas maiores ou para organizações com várias equipes.
\end{itemize}

\section{Considerações Finais}

Este capítulo apresentou a metodologia da pesquisa, classificada como aplicada,
quali-quantitativa e experimental. As duas versões da aplicação implementam os mesmos
requisitos com as mesmas tecnologias, e a sua equivalência funcional é verificada por testes
automatizados. A comparação é realizada em um ambiente conteinerizado e reproduzível, com o
mesmo orçamento de recursos para as duas arquiteturas, com a simulação do serviço de pagamento
e com cenários de carga que representam a navegação e a compra em uma loja virtual. Os
critérios de desempenho, escalabilidade, manutenibilidade, complexidade de desenvolvimento e
implantação são avaliados por meio de métricas objetivas, analisadas estatisticamente a partir
de cinco repetições de cada teste. O próximo capítulo descreve o desenvolvimento das duas
versões da aplicação.
